\nuevaReceta{Migas de Harina}{60 min}{receta:migas_harina}

    \begin{minipage}[t]{0.35\textwidth}
        \section*{Ingredientes}
        \begin{itemize}[leftmargin=*]
            \item 500 g de harina de trigo
            \item 500 ml de agua
            \item 5 Dientes de ajo
            \item 150 g de panceta
            \item 150 g de chorizo
            \item 2 Pimientos verdes
            \item 100 ml de aceite de oliva
            \item Sal
            \item Boquerones fritos para acompañar
            \item Aceitunas
        \end{itemize}
    \end{minipage}
    \hfill
    \begin{minipage}[t]{0.6\textwidth}
        \section*{Preparación}
        \begin{enumerate}
            \item \textbf{Freír los acompañamientos:} Cortar la panceta, el chorizo y los pimientos. Freírlos por separado en una sartén grande con el aceite y reservar.
            \item \textbf{Dorar los ajos:} En el mismo aceite, añadir los dientes de ajo enteros y ligeramente machacados. Cocinarlos hasta que estén dorados y reservarlos.
            \item \textbf{Añadir el agua:} Retirar la sartén del fuego y dejar que el aceite pierda un poco de temperatura. Verter el agua con cuidado, añadir sal y volver a calentar hasta que hierva.
            \item \textbf{Incorporar la harina:} Añadir toda la harina poco a poco mientras se mezcla para evitar que se formen grumos grandes.
            \item \textbf{Cocinar las migas:} Trabajar la masa a fuego medio con una rasera, cortándola, levantándola y removiéndola continuamente para separarla en migas pequeñas.
            \item \textbf{Dar el punto:} Continuar durante 25-30 minutos, hasta que las migas queden sueltas, cocinadas y ligeramente doradas.
            \item \textbf{Terminar:} Añadir la panceta, el chorizo y los pimientos reservados. Mezclar bien y servir inmediatamente con los boquerones fritos y las aceitunas.
        \end{enumerate}
    \end{minipage}

    \vspace{1em}
    \begin{tcolorbox}[colback=orange!5, colframe=orange!50, title=Consejo, size=small]
        La paciencia es fundamental: corta y remueve la masa sin parar hasta que se formen las migas. Si quedan demasiado secas, salpica unas gotas de agua; si están pegajosas, continúa cocinándolas. Los boquerones y las aceitunas aportan el contraste salado perfecto.
    \end{tcolorbox}
\end{receta}
